\documentclass[12pt]{article}
\usepackage[utf8]{inputenc}
\usepackage{ctex} % 支持中文
\usepackage{titlesec}
\usepackage{geometry}
\usepackage{hyperref} % 生成超链接目录
\usepackage{fancyhdr} % 自定义页眉页脚
\usepackage{amsmath} % 数学公式支持（必须加载以使用 \boxed）
\usepackage{booktabs} % 表格美化
\usepackage{float} % 支持 [H] 固定表格位置

% ========== 关键修复点：设置足够的页眉高度 ==========
\setlength{\headheight}{14.5pt}

% 页面设置
\geometry{a4paper, left=2.5cm, right=2.5cm, top=2.5cm, bottom=2.5cm}
\setlength{\parindent}{0pt} % 取消段落首行缩进
\setlength{\parskip}{1ex} % 段落间空一行

% 标题格式
\titleformat{\section}{\Large\bfseries}{\thesection}{1em}{}
\titleformat{\subsection}{\large\bfseries}{\thesubsection}{1em}{}
\titleformat{\subsubsection}{\normalsize\bfseries}{\thesubsubsection}{1em}{}

% 页眉页脚
\pagestyle{fancy}
\fancyhf{}
\rhead{模拟电子技术基础}
\lhead{期末总结}
\rfoot{第 \thepage 页}
\cfoot{}

% 自定义命令，避免报错
\newcommand{\Vth}{V_{\text{th}}}
\newcommand{\Vd}{V_D}
\newcommand{\Id}{I_D}
\newcommand{\Rd}{r_d}
\newcommand{\Vt}{V_T}
\newcommand{\Icbo}{I_{CBO}} % 集电结反向饱和电流
\newcommand{\Iceo}{I_{CEO}} % 穿透电流

\title{模拟电子技术基础期末总结}
\author{}
\date{}

\begin{document}

% 标题页
\maketitle
\thispagestyle{empty} % 标题页不显示页码
\newpage

% 目录
\tableofcontents
\thispagestyle{fancy}
\newpage

\section{常用半导体器件}
\subsection{半导体二极管}

\subsubsection{1. 二极管的单向导电性}
半导体二极管的核心特性是\textbf{单向导电性}，即电流仅能从二极管的\textbf{阳极（P区）}流向\textbf{阴极（N区）}，反向时二极管处于截止状态。

\textbf{导通条件}：二极管阳极电位 $V_A$ 需比阴极电位 $V_K$ 高出门槛电压（硅管典型值为 $0.7\,\text{V}$，锗管典型值为 $0.3\,\text{V}$），即满足：
\[
V_A - V_K \geq \Vth
\]
其中 $\Vth$ 为二极管的阈值电压（门槛电压）。

% =================== 后续章节占位 ===================
% 可按此格式继续添加三极管、场效应管、放大电路等内容

% =================== 参考文献（可选）===================
%\newpage
%\section*{参考文献}
%\begin{itemize}
%    \item  童诗白, 华成英. 模拟电子技术基础(第5版). 高等教育出版社, 2015.
%\end{itemize}

% =================== 尾页 ===================
%\newpage
%\thispagestyle{empty}
%\vfill
%\centering
%\Huge 感谢阅读！\\[2ex]
%\LARGE 模拟电子技术基础期末总结\\[2ex]
%\small 编制人: XXX\\[2ex]
%\small 完成日期: 2025年X月X日

